% ============================================================
% 纯答案册·课时14-2.6.2双曲线性质 —— 工具/答案抽册器.py 抽册生成件（勿手改；第三档＝纯答案单独成册）
% 源件（只读）：C:/提示词/工作区/M3-第2章量产0913/成卷/练习件/课时14-2.6.2双曲线性质/main.tex（md5 2c287b2f7d746715b5a2fb2d7aca82f4）
%% 口径：题面不重印；逐题＝题号(印面号同正册)＋[答案]＋[详解]，值/详解照源件逐字
%       （钉值门硬断言：册值≡正册件内值≡值台账值tex，逐字全等）。册序＝正册装配序＝印面号 1..16 升序（同序同号）；键序断言基准＝题面库冻结 manifest canonical 键序。
%% 三档导出：整体 main-true／纯题 main-false（在产）｜本册＝纯答案册（本件）
%% 双档：ansbook-true.tex＝详解印本；ansbook-false.tex＝纯值速查本（详解不印）
% 编译：xelatex <壳> ×2，cwd＝本目录；sty＝qp-m3.sty 本地挂载（md5 7c3930362be8a0a2bdf21bbf8ac16573，锁校验）
% ============================================================
\documentclass[fontset=none]{ctexart}
\usepackage{qp-m3}
% —— 印面逐字符号路由补钉（承源件；− 走 TNR、▱ NSC 子块；禁改 sty 保 md5） ——
\xeCJKDeclareCharClass{Default}{"2212}%
\setCJKfamilyfont{nscblk}[Path=C:/提示词/工作区/字替对照-0909/variantF/fonts/,BoldFont=NSC-w600.ttf]{NSC-w500.ttf}%
\xeCJKDeclareSubCJKBlock{nscblk}{"25B1}%
\newcommand{\pxparallelogram}{{\CJKfamily{nscblk}▱}}%
% —— 断行微调（承母版实测档，三零门承重；\emergencystretch 不另设） ——
\xeCJKsetup{CJKglue={\hskip 0.10em plus 0.08em minus 0.05em}}%
% —— 纯答案册全线括线（长详解可跨栏断，承练习件全线括线口径；灰底盒不可跨栏断） ——
\ansblockgrayfalse
% —— 双档开关（false 档＝纯值速查：答案印、详解不印；件面开关，不动 sty 本体） ——
\newif\ifansbookdetail
\ifdefined\ansbookdetail
  \ifnum\ansbookdetail=0 \ansbookdetailfalse\else\ansbookdetailtrue\fi
\else
  \ansbookdetailtrue
\fi
% —— 页脚件名（承源件章名；件名＝<件型>答案册） ——
\renewcommand{\qpzhangming}{第二章\quad 平面解析几何}%
\renewcommand{\qpjianming}{练习件答案册}%

\begin{document}
\keshi{第14课时\quad 双曲线的几何性质}
\par\glueguard{1}\addvspace{2pt}\noindent\biaoqian{【纯答案册】}{\fangsong 题面不重印\quad 印面号与正册同号同序}\par\addvspace{2pt}

\begin{multicols}{2}
\emergencystretch=1em

\begin{ansblock}[2章-练-课时14-01]
% ans:2章-练-课时14-01
\ansitem{1}{(1) 实轴长8√2，虚轴长4，顶点(−4√2,0)和(4√2,0)；(2) 实轴长6，虚轴长18，顶点(−3,0)和(3,0)；(3) 实轴长4，虚轴长4，顶点(0,−2)和(0,2)；(4) 实轴长10，虚轴长14，顶点(0,−5)和(0,5)。}
\ifansbookdetail
\ansnote{详解}{(1)化\(x^{2}/32-y^{2}/4=1\)，焦点在\(x\)轴，\(a^{2}=32\)、\(b^{2}=4\)，\(a=4\sqrt{2}\)、\(b=2\)：实轴长\(2a=8\sqrt{2}\)，虚轴长\(2b=4\)，顶点\((-4\sqrt{2},0)\)和\((4\sqrt{2},0)\)．\par\noindent (2)化\(x^{2}/9-y^{2}/81=1\)，焦点在\(x\)轴，\(a=3\)、\(b=9\)：实轴长6，虚轴长18，顶点\((-3,0)\)和\((3,0)\)．\par\noindent (3)化\(y^{2}/4-x^{2}/4=1\)，焦点在\(y\)轴，\(a=2\)、\(b=2\)：实轴长4，虚轴长4，顶点\((0,-2)\)和\((0,2)\)．\par\noindent (4)化\(y^{2}/25-x^{2}/49=1\)，焦点在\(y\)轴，\(a=5\)、\(b=7\)：实轴长10，虚轴长14，顶点\((0,-5)\)和\((0,5)\)．}
\fi
\end{ansblock}

\begin{ansblock}[2章-练-课时14-02]
% ans:2章-练-课时14-02
\ansitem{2}{实轴长18；虚轴长6；焦点(0,±3√10)；离心率 e=√10/3；渐近线方程 y=±3x。}
\ifansbookdetail
\ansnote{详解}{化\(-9x^{2}+y^{2}=81\)为\(y^{2}/81-x^{2}/9=1\)，焦点在\(y\)轴，\(a^{2}=81\)、\(b^{2}=9\)，\(a=9\)、\(b=3\)，\(c^{2}=a^{2}+b^{2}=90\)，\(c=\sqrt{90}=3\sqrt{10}\)．\par\noindent 于是：实轴长\(2a=18\)；虚轴长\(2b=6\)；焦点\((0,\pm 3\sqrt{10})\)；离心率\(e=c/a=3\sqrt{10}/9=\sqrt{10}/3\)；渐近线\(y=\pm(a/b)x=\pm 3x\)．\par\noindent 验算：\(e=\sqrt{1+(a/b)^{2}}=\sqrt{10}/3\)合；\(2c=6\sqrt{10}\approx 18.97>2a=18\)合；点\((0,9)\)代入方程\(81/81=1\)合．}
\fi
\end{ansblock}

\begin{ansblock}[2章-练-课时14-03]
% ans:2章-练-课时14-03
\ansitem{3}{B}
\ifansbookdetail
\ansnote{详解}{由\(x^{2}-y^{2}/4=1\)知焦点在\(x\)轴，\(a=1\)、\(b=2\)，渐近线\(y=\pm(b/a)x=\pm 2x\)，故选B．}
\fi
\end{ansblock}

\begin{ansblock}[2章-练-课时14-04]
% ans:2章-练-课时14-04
\ansitem{4}{3}
\ifansbookdetail
\ansnote{详解}{\(c^{2}=16+9=25\)，\(c=5\)，焦点\((\pm 5,0)\)；渐近线方程\(3x\pm 4y=0\)．取\(F(5,0)\)与直线\(3x-4y=0\)：\(d=|3\times 5-4\times 0|/\sqrt{3^{2}+4^{2}}=15/5=3\)．由对称性，任一焦点到任一条渐近线的距离均为3（恰为半虚轴长\(b=3\)）．}
\fi
\end{ansblock}

\begin{ansblock}[2章-练-课时14-05]
% ans:2章-练-课时14-05
\ansitem{5}{2}
\ifansbookdetail
\ansnote{详解}{焦点在\(x\)轴，渐近线\(y=\pm(b/a)x\)，故\(b/a=\sqrt{3}\)．\(e=c/a=\sqrt{1+(b/a)^{2}}=\sqrt{1+3}=2\)．}
\fi
\end{ansblock}

\begin{ansblock}[2章-练-课时14-06]
% ans:2章-练-课时14-06
\ansitem{6}{5/4（焦点在x轴）或 5/3（焦点在y轴）}
\ifansbookdetail
\ansnote{详解}{焦点在\(x\)轴时，渐近线\(y=\pm(b/a)x\)，\(b/a=3/4\)，\(e=\sqrt{1+(b/a)^{2}}=\sqrt{1+9/16}=\sqrt{25/16}=5/4\)；焦点在\(y\)轴时，渐近线\(y=\pm(a/b)x\)，\(a/b=3/4\)，\(b/a=4/3\)，\(e=\sqrt{1+16/9}=\sqrt{25/9}=5/3\)．题面未限焦点位置，两解并立．}
\fi
\end{ansblock}

\begin{ansblock}[2章-练-课时14-07]
% ans:2章-练-课时14-07
\ansitem{7}{D}
\ifansbookdetail
\ansnote{详解}{实轴长\(8\)得\(a=4\)；焦点在\(y\)轴，渐近线\(y=\pm(a/b)x=\pm(4/b)x=(4/3)x\)，得\(b=3\)．故\(C\)：\(y^{2}/16-x^{2}/9=1\)，选D．}
\fi
\end{ansblock}

\begin{ansblock}[2章-练-课时14-08]
% ans:2章-练-课时14-08
\ansitem{8}{(1) 实轴长8/3，虚轴长4，焦点(±2√13/3,0)，顶点(±4/3,0)，渐近线 y=±(3/2)x；(2)（m>0）实轴长4，虚轴长2√m，焦点(±√(m+4),0)，顶点(±2,0)，渐近线 y=±(√m/2)x。}
\ifansbookdetail
\ansnote{详解}{(1)\(a^{2}=16/9\)、\(b^{2}=4\)，\(a=4/3\)、\(b=2\)；\(c^{2}=16/9+4=52/9\)，\(c=\sqrt{52/9}=2\sqrt{13}/3\)：实轴长\(8/3\)，虚轴长4，焦点\((\pm 2\sqrt{13}/3,0)\)，顶点\((\pm 4/3,0)\)，渐近线\(y=\pm(b/a)x=\pm(3/2)x\)．\par\noindent (2)双曲线须\(m>0\)：\(a=2\)、\(b=\sqrt{m}\)、\(c^{2}=4+m\)，\(c=\sqrt{m+4}\)：实轴长4，虚轴长\(2\sqrt{m}\)，焦点\((\pm\sqrt{m+4},0)\)，顶点\((\pm 2,0)\)，渐近线\(y=\pm(\sqrt{m}/2)x\)．（\(m<0\)时方程表椭圆，与题设矛盾，仅\(m>0\)一解．）}
\fi
\end{ansblock}

\begin{ansblock}[2章-练-课时14-09]
% ans:2章-练-课时14-09
\ansitem{9}{x²/3−y²=1}
\ifansbookdetail
\ansnote{详解}{由渐近线得\(b/a=\sqrt{3}/3\)，即\(a=\sqrt{3}b\)．直线\(l\)过\(P(0,b)\)、\(Q(a,0)\)，方程\(bx+ay-ab=0\)，原点到\(l\)的距离\(ab/\sqrt{a^{2}+b^{2}}=\sqrt{3}/2\)，两边平方得\(4a^{2}b^{2}=3(a^{2}+b^{2})\)．\par\noindent 联立\(a=\sqrt{3}b\)：\(4\times 3b^{4}=3(3b^{2}+b^{2})\)，即\(12b^{4}=12b^{2}\)，\(b^{2}=1\)（\(b>0\)），\(a^{2}=3\)．故双曲线方程为\(x^{2}/3-y^{2}=1\)．验算：\(a=\sqrt{3}\)、\(b=1\)时\(ab/\sqrt{a^{2}+b^{2}}=\sqrt{3}/2\)合，渐近线\(y=\pm x/\sqrt{3}\)合．}
\fi
\end{ansblock}

\begin{ansblock}[2章-练-课时14-10]
% ans:2章-练-课时14-10
\ansitem{10}{(1) x²/9−y²/25=1 与 y²/25−x²/9=1（答案不唯一，凡渐近线为 5x±3y=0 的标准方程均可）；(2) y²/(56/9)−x²/(56/25)=1（即 9y²−25x²=56）。}
\ifansbookdetail
\ansnote{详解}{(1)两渐近线即\(y=\pm(5/3)x\)．焦点在\(x\)轴：\(b/a=5/3\)，取\(a=3\)、\(b=5\)得\(x^{2}/9-y^{2}/25=1\)；焦点在\(y\)轴：\(a/b=5/3\)，取\(a=5\)、\(b=3\)得\(y^{2}/25-x^{2}/9=1\)．\par\noindent (2)以\(l_1\)、\(l_2\)为渐近线的双曲线可设为\(x^{2}/9-y^{2}/25=\lambda\)（\(\lambda\neq 0\)）．代入\(M(1,3)\)：\(\lambda=1/9-9/25=(25-81)/225=-56/225<0\)，故焦点在\(y\)轴，方程即\(y^{2}/25-x^{2}/9=56/225\)，化标准形\(y^{2}/(56/9)-x^{2}/(56/25)=1\)．验算：\(9y^{2}-25x^{2}=56\)代入\(M(1,3)\)：\(81-25=56\)合．提醒：渐近线同为\(5x\pm 3y=0\)的双曲线焦点可在\(x\)轴或\(y\)轴，由所过点的坐标定号（本题\(1/9-9/25<0\)定\(y\)轴）．}
\fi
\end{ansblock}

\begin{ansblock}[2章-练-课时14-11]
% ans:2章-练-课时14-11
\ansitem{11}{B}
\ifansbookdetail
\ansnote{详解}{\(F(-c,0)\)，\(A(a,0)\)，取\(B(0,b)\)，直线\(AB\)：\(bx+ay-ab=0\)．\(F\)到\(AB\)的距离\(d=|-bc-ab|/\sqrt{a^{2}+b^{2}}=b(c+a)/c\)（分子提出\(b\)，分母用\(c^{2}=a^{2}+b^{2}\)）．\par\noindent 由\(d=(5/4)b\)得\(c+a=(5/4)c\)，即\(4a=c\)；\(b^{2}=c^{2}-a^{2}=15a^{2}\)，\(b/a=\sqrt{15}\)，渐近线\(y=\pm\sqrt{15}x\)，故选B．验算：\(c=4a\)时\(d=b\cdot 5a/(4a)=(5/4)b\)合．}
\fi
\end{ansblock}

\begin{ansblock}[2章-练-课时14-12]
% ans:2章-练-课时14-12
\ansitem{12}{A}
\ifansbookdetail
\ansnote{详解}{设\(M(x_0,y_0)\)，则\(N(-x_0,-y_0)\)，\(P(x,y)\)（\(x\neq\pm x_0\)）．\(k_1k_2=\dfrac{y-y_0}{x-x_0}\cdot\dfrac{y+y_0}{x+x_0}=\dfrac{y^{2}-y_0^{2}}{x^{2}-x_0^{2}}\)（坐标点差法，不引名目）．\(M\)、\(P\)均在双曲线上：\(y^{2}=b^{2}(x^{2}/a^{2}-1)\)，\(y_0^{2}=b^{2}(x_0^{2}/a^{2}-1)\)，相减得\(y^{2}-y_0^{2}=(b^{2}/a^{2})(x^{2}-x_0^{2})\)，故\(k_1k_2=b^{2}/a^{2}\)（定值）．\par\noindent 于是\(|k_1|+|k_2|\geqslant 2\sqrt{|k_1k_2|}=2b/a\)（当\(|k_1|=|k_2|=b/a\)时取等，\(P\)可取到），由最小值2得\(b/a=1\)，\(e=\sqrt{1+(b/a)^{2}}=\sqrt{2}\)，故选A．}
\fi
\end{ansblock}

\begin{ansblock}[2章-练-课时14-13]
% ans:2章-练-课时14-13
\ansitem{13}{1}
\ifansbookdetail
\ansnote{详解}{\(c^{2}=16+9=25\)，\(c=5\)，右焦点\(F(5,0)\)．设\(P(x,y)\)，\(|PF|^{2}=(x-5)^{2}+y^{2}\)．由\(y^{2}=9(x^{2}/16-1)\)（\(|x|\geqslant 4\)）：\(|PF|^{2}=x^{2}-10x+25+(9/16)x^{2}-9=(25/16)x^{2}-10x+16\)．对称轴\(x=10/(2\times 25/16)=16/5<4\)，故在\(|x|\geqslant 4\)上于\(x=4\)取最小：\(|PF|^{2}=(25/16)\times 16-40+16=1\)，\(|PF|_{\min}=1\)（即右顶点处取得，同支焦半径最小值\(c-a=5-4=1\)合）．}
\fi
\end{ansblock}

\begin{ansblock}[2章-练-课时14-14]
% ans:2章-练-课时14-14
\ansitem{14}{证明见解析。}
\ifansbookdetail
\ansnote{详解}{证明：不妨先设焦点在\(x\)轴：双曲线\(x^{2}/a^{2}-y^{2}/b^{2}=1\)（\(a>0\)，\(b>0\)），焦点\(F(c,0)\)，渐近线\(y=\pm(b/a)x\)，即\(bx\mp ay=0\)．焦点到渐近线的距离\(d=|bc-a\cdot 0|/\sqrt{a^{2}+b^{2}}=bc/c=b\)（半虚轴长）．\par\noindent 焦点在\(y\)轴时（\(y^{2}/a^{2}-x^{2}/b^{2}=1\)，焦点\(F(0,c)\)，渐近线\(y=\pm(a/b)x\)即\(ay\mp bx=0\)）：\(d=|a\cdot 0-bc|/\sqrt{a^{2}+b^{2}}=bc/c=b\)，等式仍成立．证毕．}
\fi
\end{ansblock}

\begin{ansblock}[2章-练-课时14-15]
% ans:2章-练-课时14-15
\ansitem{15}{(1) e=√5/2；(2) 证明见解析，定点坐标为 (−10/3,0)。}
\ifansbookdetail
\ansnote{详解}{(1)\(a=2\)，\(b=1\)，\(c=\sqrt{4+1}=\sqrt{5}\)，\(e=c/a=\sqrt{5}/2\)。\par\noindent (2)设\(A(x_1,y_1)\)，\(B(x_2,y_2)\)。联立\(y=kx+m\)与\(x^{2}/4-y^{2}=1\)：\((1-4k^{2})x^{2}-8mkx-4(m^{2}+1)=0\)，其中\(1-4k^{2}\neq 0\)（否则\(l\)与渐近线平行，不合两交点），\(\Delta=64m^{2}k^{2}+16(1-4k^{2})(m^{2}+1)>0\)；\(x_1+x_2=8mk/(1-4k^{2})\)，\(x_1x_2=-4(m^{2}+1)/(1-4k^{2})\)；\(y_1y_2=(kx_1+m)(kx_2+m)=k^{2}x_1x_2+mk(x_1+x_2)+m^{2}=(m^{2}-4k^{2})/(1-4k^{2})\)。\par\noindent 以\(AB\)为直径的圆过\(D(-2,0)\Leftrightarrow\overrightarrow{DA}\cdot\overrightarrow{DB}=0\Leftrightarrow(x_1+2)(x_2+2)+y_1y_2=0\)，即\(x_1x_2+2(x_1+x_2)+4+y_1y_2=0\)。通分（分母\(1-4k^{2}\)）：\([m^{2}-4k^{2}]+[-4(m^{2}+1)]+[16mk]+[4(1-4k^{2})]=-3m^{2}+16mk-20k^{2}=0\)，即\(3m^{2}-16mk+20k^{2}=0\)，解得\(m=2k\)或\(m=(10/3)k\)。\par\noindent 当\(m=2k\)时，\(l\)：\(y=k(x+2)\)过\(D(-2,0)\)，此时\(A\)、\(B\)之一即\(D\)，与“\(A\)、\(B\)均异于左、右顶点”矛盾，舍去；当\(m=(10/3)k\)时（\(k\neq 0\)），\(l\)：\(y=k(x+10/3)\)恒过定点\((-10/3,0)\)，且不过\(D\)，符合题意。故直线\(l\)过定点，定点坐标为\((-10/3,0)\)。验算：\(3m^{2}-16mk+20k^{2}=(3m-10k)(m-2k)\)合；定点\((-10/3,0)\)在左顶点\(D(-2,0)\)左侧，直线族绕该点转动时与双曲线恒有两交点（\(k\)充分小时），满足\(\Delta>0\)合。}
\fi
\end{ansblock}

\begin{ansblock}[2章-练-课时14-16]
% ans:2章-练-课时14-16
\ansitem{16}{(1) y=±√3x；(2) 存在，定直线方程 x=1/2。}
\ifansbookdetail
\ansnote{详解}{(1)\(|OF_1|=|OF_2|=c=2\)，\(|OP|=|OF_1|\)说明\(\triangle PF_1F_2\)中边\(F_1F_2\)上的中线等于该边一半，故\(\angle F_1PF_2=90^{\circ}\)，即\(PF_1\perp PF_2\)。\par\noindent \(P\)在右支：\(|PF_1|-|PF_2|=2a\)；面积：\(|PF_1|\cdot|PF_2|=2S=6\)。\par\noindent \(|PF_1|^{2}+|PF_2|^{2}=|F_1F_2|^{2}=(2c)^{2}=16=(|PF_1|-|PF_2|)^{2}+2|PF_1||PF_2|=4a^{2}+12\)，故\(4a^{2}=4\)，\(a=1\)。\par\noindent \(b^{2}=c^{2}-a^{2}=4-1=3\)，渐近线\(y=\pm(b/a)x=\pm\sqrt{3}x\)。\par\noindent (2)双曲线\(C\)：\(x^{2}-y^{2}/3=1\)。\par\noindent （i）\(l\)斜率不存在：\(l\)：\(x=2\)，\(M(2,3)\)，\(N(2,-3)\)。直线\(A_1M\)：\(y=x+1\)；直线\(A_2N\)：\(y=-3x+3\)。联立得\(Q(1/2,3/2)\)，横坐标\(1/2\)。\par\noindent （ii）\(l\)斜率存在：\(l\)过\(F_2(2,0)\)且不与渐近线平行，设\(l\)：\(y=k(x-2)\)（\(k\neq 0\)，\(k\neq\pm\sqrt{3}\)），\(M(x_1,y_1)\)，\(N(x_2,y_2)\)。联立得\((3-k^{2})x^{2}+4k^{2}x-4k^{2}-3=0\)，\(x_1+x_2=4k^{2}/(k^{2}-3)\)，\(x_1x_2=(4k^{2}+3)/(k^{2}-3)\)。\par\noindent 直线\(A_1M\)：\(y=[y_1/(x_1+1)](x+1)\)；直线\(A_2N\)：\(y=[y_2/(x_2-1)](x-1)\)。交点\(Q(x,y)\)满足\((x+1)/(x-1)=y_2(x_1+1)/[y_1(x_2-1)]\)（\(x\neq 1\)）。两边平方，以\(y_1^{2}=3(x_1^{2}-1)\)、\(y_2^{2}=3(x_2^{2}-1)\)代入：\par\noindent \(\left[(x+1)/(x-1)\right]^{2}=\dfrac{\left[3(x_2^{2}-1)(x_1+1)^{2}\right]}{\left[3(x_1^{2}-1)(x_2-1)^{2}\right]}\)\par\noindent \(=\left[(x_2+1)(x_1+1)\right]/\left[(x_1-1)(x_2-1)\right]\)\par\noindent \(=\left[x_1x_2+(x_1+x_2)+1\right]/\left[x_1x_2-(x_1+x_2)+1\right]\)\par\noindent \(=\left[(4k^{2}+3)+4k^{2}+(k^{2}-3)\right]/\left[(4k^{2}+3)-4k^{2}+(k^{2}-3)\right]\)\par\noindent \(=9k^{2}/k^{2}=9\)。\par\noindent 故\((x+1)/(x-1)=\pm 3\)，得\(x=1/2\)或\(x=2\)。\par\noindent 取\(x=2\)需\(y_2(x_1+1)/[y_1(x_2-1)]=+3>0\)；但\(M\)、\(N\)是过右焦点的弦的两端点且均非顶点：\(0<k^{2}<3\)时\(M\)、\(N\)分居两支（\(x_1x_2<0\)），\(x_2-1<0\)而\(y_1y_2>0\)；\(k^{2}>3\)时\(M\)、\(N\)同在右支且分居\(x\)轴两侧，\(y_1y_2<0\)而\((x_1+1)(x_2-1)>0\)——两种情形该式恒为负，只能取\(-3\)。故\(x=2\)为平方增根，舍去，\(x=1/2\)。\par\noindent 综上，交点\(Q\)恒在定直线\(x=1/2\)上。\par\noindent 验算：（i）的\(Q(1/2,3/2)\)落在\(x=1/2\)上合；\par\noindent （ii）韦达分子分母化简：\par\noindent \(4k^{2}+3+4k^{2}+k^{2}-3=9k^{2}\)、\(4k^{2}+3-4k^{2}+k^{2}-3=k^{2}\)，比值9合。}
\fi
\end{ansblock}

% ---- 尾块（\tailfill[30mm] 定高参——钉档 §三.3：无参在平衡栏呈「内容尾随框」，贴栏底须定高参；实测无参致末页平衡 Overfull\vbox 12.99pt，定高 30mm 三零） ----
\tailfill[30mm]

\end{multicols}
\end{document}
